\chapter{Literature Review}
\label{ch:literature}

\section{Introduction}

This chapter looks at the research and industry work that shaped Echo's
design. It has three parts. Section~\ref{sec:lit-password} looks at why
passwords keep failing, even after decades of security advice.
Section~\ref{sec:lit-passwordless} compares the passwordless methods used
today, using the comparison approach set out by Bonneau and colleagues
\cite{bonneau2012quest}, so the trade-offs are shown clearly rather than
just described in general terms. Section~\ref{sec:lit-acoustic} looks at
research on sending data through sound and pairing devices, which is what
Echo's sound channel is built on. Section~\ref{sec:lit-position} brings
these three parts together to explain exactly what Echo adds that is new.

\section{Why Password Authentication Fails}
\label{sec:lit-password}

A password asks a user to invent, remember, and correctly recall a made-up
secret for every account they own, without ever writing it down somewhere
unsafe, reusing it, or being fooled by a fake copy of the login page asking
for it. NIST, the US government's standards body, changed its official
password guidance for exactly this reason: it moved away from forcing
users to change passwords often and follow strict complexity rules,
because those older rules pushed people toward predictable passwords that
were actually easier, not harder, to guess \cite{nist80063b}. Instead,
NIST now recommends checking new passwords against lists of already-leaked
passwords, allowing long plain-English phrases instead of complex strings,
and relying on a second login factor for anything important. In short,
even the standards body admits that a password alone is no longer good
enough to trust.

The cost of this weakness shows up clearly in real breach numbers. The
2024 Verizon Data Breach Investigations Report found that stolen, reused,
or guessed passwords were the leading way attackers broke into systems
across 10,626 confirmed breaches, with phishing close behind and human
error playing a part in over two-thirds of all cases
\cite{verizon2024dbir}. OWASP, a well-known security organisation, lists
in its Authentication Cheat Sheet everything a developer must add just to
keep a normal password login reasonably safe: multi-factor login,
protection against stolen-password reuse, limits on repeated attempts, safe
cookie settings, and checks against leaked-password lists
\cite{owaspauthn}. Every one of these is a patch added on top of a weak
idea, rather than a fix to the idea itself. That is the real reason the
industry has spent the last ten years building alternatives.

\section{Passwordless Authentication: A Comparative Review}
\label{sec:lit-passwordless}

Bonneau and colleagues \cite{bonneau2012quest} built a framework for
comparing password replacements on ease of use, ease of adoption, and
security, and their main finding still holds true more than ten years
later: not one method they studied kept every benefit that ordinary
passwords already had. In other words, every passwordless method below
trades away something in exchange for something else.

\subsection{One-Time Codes and Magic Links}

Text message and email one-time codes fix the password reuse problem
directly, since a stolen password database cannot be reused to log in with
a one-time code. But the code still has to pass through something the
user reads and retypes, so it can still be phished: a fake login page can
simply ask for the code and pass it on to the real site straight away.
Text-message codes also depend on the mobile network staying secure, and
are open to a well-known attack called SIM-swapping, which is exactly why
NIST's guidelines discourage using text messages as a login factor for
anything sensitive \cite{nist80063b}. Magic links remove the retyping step
but add a separate channel (email) that the user has to switch to, so the
login is now only as safe as that email account.

\subsection{QR-Code Login}

QR-code login, the pattern made popular by WhatsApp Web and used by many
workplace login systems, shows a code on the login screen that an
already-logged-in phone scans to approve the session. This removes typing
completely and ties the approval to a device the user already trusts, but
it means aiming a camera at a screen, which is awkward in some situations.
It also has its own phishing risk: if an attacker can copy the real site's
QR code onto their own fake page, the code itself does not prove which
physical device is actually asking for it.

\subsection{FIDO2 / WebAuthn Passkeys}

The FIDO2 and WebAuthn standards, built jointly by the FIDO Alliance and
the W3C, are where the industry is currently heading for passwordless
login, and they are the technology behind "sign in with a passkey" on
Apple, Google, and Microsoft devices since 2022
\cite{fidoalliance, w3cwebauthn}. A passkey is a pair of matching digital
keys created on the user's own device; the private half never leaves that
device (or the company's secure syncing system), and the browser ties
every login attempt to the exact website asking for it. This makes
passkeys very hard to phish: a fake copy of a login page hosted on the
wrong website simply cannot get a valid login approval, no matter how
convincing it looks. This is the strongest of the methods covered here,
and Echo borrows its main idea directly: a key pair tied to one device,
where the private half never leaves that device.

Passkeys are not perfectly smooth to use, though. The user still has to
notice and respond to a pop-up, choose the right device or account when
more than one is available, and, when logging in on a new device, often
has to scan a QR code anyway, similar to the method described above.
Syncing passkeys across a person's Apple or Google account is convenient,
but it locks the user into that company's system, and the way passkeys
look and behave still differs noticeably between browsers and phones,
which is part of why they have not yet fully replaced passwords as a
fallback option on most websites.

\subsection{Biometric and Proximity-Based Approaches}

Fingerprint and face recognition on a phone are usually used today as a
lock on an existing key, for example, unlocking a passkey, rather than
being sent to a remote server directly, since fingerprint data is not
something that should ever be sent over the internet or stored elsewhere.
Bluetooth- and NFC-based login (used in some workplace systems) is closer
to what Echo attempts: it logs a user in because a trusted device is
nearby, without needing the user to unlock or tap that device every time
after an initial setup. The problem is cost and inconsistency: Bluetooth
pairing and distance detection behave differently across phone brands and
increasingly require the user to grant special permissions, and NFC needs
the two devices to be brought close enough to physically touch, which
removes the "just walk up and be recognised" feeling this type of login is
supposed to give.

\section{Acoustic Data Transmission and Device Pairing}
\label{sec:lit-acoustic}

Using sound to send data between everyday devices is an old idea, older
than the smartphone. Madhavapeddy and colleagues
\cite{madhavapeddy2005audio} showed that ordinary audio tones could carry
useful data between two mobile phones during an active phone call, sending
around 20 bits per second with a very low error rate over 3 metres, using
DTMF-style tones (the same kind used for phone keypads). Their work showed
that sound is a genuinely usable wireless channel, not just a novelty,
because it needs no Bluetooth pairing, no Wi-Fi connection, and no direct
line of sight. That same idea, that ordinary speakers and microphones can
carry data with no setup step at login time, is exactly what makes sound
useful as proof that a phone is physically nearby.

Security researchers have also studied sound channels from the opposite
direction, as a way data could leak out without permission. Hanspach and
Goetz \cite{hanspach2013covert} built a network of computers that could
pass data to each other, hop by hop, using nothing but their normal
speakers and microphones at a near-silent pitch, in the very same
18 to 20\,kHz range that Echo uses, and suggested ways to defend against it,
such as filtering out high-pitched sound in hardware and watching for
unexpected microphone use. Their work is a useful warning for Echo's own
threat model (Chapter~\ref{ch:design}): the same qualities that make a
near-silent sound channel convenient for a normal login, it passes through
some walls, needs no pairing, and cannot be heard, also make it useful for
a malicious data leak on a device that has already been broken into. A
security-minded system should be open about when and why it is using the
microphone or speaker.

The \texttt{ggwave} library \cite{ggwave}, which Echo uses to send and
receive sound, works by playing several tones at once to represent data
(a method called Frequency-Shift Keying) with built-in error correction.
It can send roughly 8 to 16 bytes per second depending on the settings used,
which is more than enough for a short one-time code and phone identifier.
It offers an \texttt{ULTRASOUND\_NORMAL} setting in the 18 to 20\,kHz range,
plus a normal audible fallback for hardware that cannot reproduce those
high frequencies cleanly, and Echo uses both directly.

\section{Positioning Echo}
\label{sec:lit-position}

Looking at all this research together shows a gap worth filling. The
individual pieces Echo needs, a key pair tied to one device (from
WebAuthn), a no-setup-needed nearby-device channel (from audio networking
research), and an honest understanding of that channel's own risks (from
research on sound as a data leak), are each well studied on their own, but
rarely put together into one working, understandable website login system
outside of the big tech companies' own passkey systems. Echo does not
invent a new cryptographic method; it uses the exact same ECDSA signing
that WebAuthn uses. What Echo adds is a complete, understandable login
process that replaces the device-discovery step passkeys normally handle
behind the scenes with a simple sound-based handshake, judged honestly
against both what it proves cryptographically
(Chapter~\ref{ch:design}, Section~\ref{sec:testing}) and how reliable the
sound channel actually is in the real world
(Chapter~\ref{ch:design}, Section~\ref{sec:evaluation};
Chapter~\ref{ch:conclusion}).
